\chapter{Materials and textures}
\label{ch:materials}

\section{Collecting the materials}
The ten example materials come from two libraries of physically based materials released under
the CC0 licence (public domain): seven from \emph{ambientCG} and three from \emph{Poly Haven}
(Appendix~\ref{app:materials}). The script \file{tools/fetch\_textures.ps1} downloads them:
\begin{itemize}
\item For ambientCG it asks the web API once for all seven assets
(\url{https://ambientcg.com/api/v2/full_json?id=...&include=downloadData}), downloads the
\code{1K-JPG} ZIP archive of each, and extracts the colour, displacement, roughness and DirectX
normal maps.
\item For Poly Haven it asks \url{https://api.polyhaven.com/files/<id>} for the address of each map
at 1K resolution, and \url{https://api.polyhaven.com/info/<id>} for the name and authors.
\item It renames every map to a standard name (\file{color.jpg}, \file{displacement.jpg},
\file{roughness.jpg}, \file{normal.jpg}) and writes a \emph{manifest}.
\end{itemize}
The manifest records each material's origin, so the program can credit its source:
\begin{lstlisting}[language={},numbers=none,basicstyle=\ttfamily\footnotesize]
{
  "id": "Bricks076C",  "name": "Bricks 076 C",  "site": "ambientCG",
  "url": "https://ambientcg.com/view?id=Bricks076C",
  "author": "Lennart Demes (ambientCG)",  "license": "CC0 1.0",
  "files": { "color": "Bricks076C/color.jpg", "displacement": "Bricks076C/displacement.jpg",
             "roughness": "Bricks076C/roughness.jpg", "normal": "Bricks076C/normal.jpg" },
  "bumpFromHeight": true
}
\end{lstlisting}
Neither library ships a separate \emph{bump} map for these materials, and that is no accident:
a bump map \emph{is} a height field, and the displacement map is the height field of the material.
The program therefore uses the displacement map as the bump map when no bump map is given, and
labels it ``Bump (from height)'' in the interface.

\section{Two copies of every map}
Each map is kept twice. The GPU copy is a texture with a complete mipmap chain, used for
rendering and for display in the interface. The CPU copy, an array of bytes, is what the probe
reads. Keeping it costs 4\,MB per $1024\times1024$ map, a small price for being able to compute
on the CPU exactly what the GPU computes.

\excerpt{mapimage}

\section{Sampling on the CPU}
\label{sec:cpusampling}
\code{MapImage::Sample} is the bilinear filter of Chapter~\ref{ch:sampling}, with wrap
addressing and the GPU's texel-centre convention (the $-0.5$):

\excerpt{sample}

\section{Loading an image into a typeless texture}
\label{sec:typeless}
\code{LoadImage} decodes the file with \code{stb\_image}, always asking for four channels (a grey
JPEG becomes grey RGB, which keeps the rest of the code uniform), and creates the texture. The
important details are in the texture description:
\begin{itemize}
\item \code{DXGI\_FORMAT\_R8G8B8A8\_TYPELESS} allows two views with different interpretations of
the same bytes: \code{UNORM} returns the stored values; \code{UNORM\_SRGB} decodes sRGB to linear
(Section~\ref{sec:srgb}). Only the colour map gets an sRGB view; data maps must not be decoded.
\item \code{MipLevels = 0} requests the full chain, and the flags \code{RENDER\_TARGET} and
\code{GENERATE\_MIPS} let the GPU compute the smaller levels with \code{GenerateMips}.
\end{itemize}

\excerpt{loadimage}

\section{Loading a material}
A material's maps are loaded the first time it is needed; the program loads all ten at start-up
so that the list can show thumbnails. Note how the bump map falls back to the displacement map:

\excerpt{ensureloaded}
